% =====================================================================
\chapter{Soluciones de los ejercicios}
\label{ap:soluciones}
% =====================================================================

Cada solución se presenta en cinco partes: datos, planteamiento, desarrollo, resultado y comprobación. En los pasos intermedios se redondea a cuatro cifras significativas salvo que se indique otra cosa. Los ejercicios de diseño razonado no tienen una única respuesta correcta; la solución propone una y justifica cada decisión.

% ---------------------------------------------------------------------
\section{Tema 1}

\begin{solucion}{eje:1-clasificar}
\datos cuatro situaciones.
\planteamiento aplicar la \cref{def:cajanegra} y la \cref{def:regret}: ¿se conoce $f$ como fórmula?, ¿hay gradiente?, ¿hay ruido?, ¿importa cada ensayo o solo el final?
\pasos
\begin{enumerate}[label=(\alph*)]
  \item Horno de pan: caja negra, ruidosa (variabilidad de la masa y de la medida de calidad), sin gradiente. Las hornadas de prueba se venden: cada ensayo malo cuesta, luego interesa el \emph{regret acumulado}.
  \item Ala con CFD: caja negra (salvo que se disponga del adjunto del código), determinista, cara. Solo importa el diseño final: \emph{regret simple}.
  \item Red con descenso de gradiente: no es caja negra; hay gradiente por diferenciación automática y evaluaciones baratas. No se aplican los métodos del módulo.
  \item Tasa de aprendizaje de un LLM: caja negra, ruidosa (semilla, orden de los datos), sin gradiente respecto al hiperparámetro y con fidelidades baratas (ejecuciones cortas). Regret simple sobre la configuración final.
\end{enumerate}
\resultado \textbf{(a), (b) y (d) son problemas de caja negra; (a) y (d) con ruido; solo (c) tiene gradiente; regret acumulado en (a) y simple en (b) y (d).}
\comprobacion la clasificación coincide con la \cref{tab:familias}: (c) pertenece a la fila de métodos de gradiente.
\end{solucion}

\begin{solucion}{eje:1-rejilla}
\datos (a) 4 niveles, 5 variables; (b) $h=0{,}05$, $d\in\{2,6\}$; (c) 125 ensayos de 2~h.
\planteamiento regla del producto y \cref{prop:rejilla}: $N\ge(2h)^{-d}$.
\pasos
\begin{enumerate}[label=\arabic*.]
  \item $4^5=1024$ ensayos.
  \item $(2h)^{-1}=1/0{,}1=10$; con $d=2$, $N\ge10^2=100$; con $d=6$, $N\ge10^6$.
  \item $125\cdot2=250$~h; $250/24=10{,}4$ días.
\end{enumerate}
\resultado \textbf{1024 ensayos; al menos 100 y $10^6$ puntos; 250 horas, unos 10,4 días de máquina sin parar.}
\comprobacion multiplicar la dimensión por 3 eleva la cota al cubo: $(10^2)^3=10^6$.
\end{solucion}

\begin{solucion}{eje:1-coste}
\datos $x_1=(230,60,19)$, $d_1=0{,}0317$, $w_1=0{,}480$~mm; $x_2=(225,65,22)$, $d_2=0{,}0306$, $w_2=0{,}421$~mm.
\planteamiento $t_c=8+t_e+0{,}05(T-200)$ y $C=(45t_c/3600+0{,}40)/(1-d)$; comprobar $w\le0{,}45$.
\pasos
\begin{enumerate}[label=\arabic*.]
  \item $t_{c,1}=8+19+1{,}5=28{,}5$~s; $45\cdot28{,}5/3600=0{,}3563$; $C_1=0{,}7563/0{,}9683=0{,}7810$~€.
  \item $t_{c,2}=8+22+1{,}25=31{,}25$~s; $45\cdot31{,}25/3600=0{,}3906$; $C_2=0{,}7906/0{,}9694=0{,}8156$~€.
  \item $w_1=0{,}480>0{,}45$: $x_1$ no es factible. $w_2=0{,}421\le0{,}45$: $x_2$ sí.
\end{enumerate}
\resultado \textbf{$C(x_1)=0{,}7810$~€ pero incumple el alabeo; $C(x_2)=0{,}8156$~€ y es factible: se prefiere $x_2$.}
\comprobacion ambos costes están entre el óptimo ($0{,}7782$) y la receta del operario ($0{,}8688$), como debe ser para ajustes intermedios.
\end{solucion}

\begin{solucion}{eje:1-max2}
\datos $Z_1,Z_2\sim\N(0,1)$ independientes; $\sigma=0{,}0126$~€.
\planteamiento $\max(a,b)=\frac12(a+b)+\frac12\abs{a-b}$ y $\E\abs{W}=s\sqrt{2/\pi}$ si $W\sim\N(0,s^2)$.
\pasos
\begin{enumerate}[label=\arabic*.]
  \item $\E[\max(Z_1,Z_2)]=\frac12\E[Z_1+Z_2]+\frac12\E\abs{Z_1-Z_2}=0+\frac12\E\abs{Z_1-Z_2}$.
  \item $Z_1-Z_2\sim\N(0,2)$, luego $\E\abs{Z_1-Z_2}=\sqrt2\sqrt{2/\pi}=2/\sqrt\pi$.
  \item $\E[\max]=\frac12\cdot2/\sqrt\pi=1/\sqrt\pi=0{,}5642$.
  \item Para la inyectora: el mejor de dos lotes iguales parece $0{,}5642\cdot0{,}0126=0{,}0071$~€ más barato.
\end{enumerate}
\resultado \textbf{$\E[\max(Z_1,Z_2)]=1/\sqrt\pi\approx0{,}5642$; sesgo de unos 0,7 céntimos por pieza.}
\comprobacion la práctica estima por Monte Carlo $0{,}561$ con $2\cdot10^5$ repeticiones, compatible con el valor exacto.
\end{solucion}

\begin{solucion}{eje:1-lote}
\datos $d=0{,}02$; desviación deseada $0{,}005$; lote alternativo $m=200$.
\planteamiento por la \cref{prop:binomial}, $\Var(\hat d)=d(1-d)/m$.
\pasos
\begin{enumerate}[label=\arabic*.]
  \item $m=d(1-d)/0{,}005^2=0{,}0196/0{,}000\,025=784$ piezas.
  \item Con $m=200$: $\sqrt{0{,}0196/200}=\sqrt{0{,}000\,098}=0{,}0099$, es decir, 0,99 puntos.
  \item Promediar $k$ lotes de 200 equivale a un lote de $200k$; $200k\ge784$ exige $k=4$ lotes ($3{,}92$ redondeado).
\end{enumerate}
\resultado \textbf{784 piezas; 0,99 puntos porcentuales con 200; hacen falta 4 lotes de 200.}
\comprobacion la precisión mejora como $1/\sqrt m$: cuadruplicar el lote (de 200 a 800) reduce la desviación a la mitad, de 0,99 a 0,49 puntos.
\end{solucion}

\begin{solucion}{eje:1-sustituto}
\datos 40~min por simulación; 400 evaluaciones directas; 60 + 10 con sustituto; 2~h de trabajo.
\planteamiento comparar tiempos totales y enumerar riesgos.
\pasos
\begin{enumerate}[label=\arabic*.]
  \item Directo: $400\cdot40=16\,000$~min $=266{,}7$~h.
  \item Sustituto: $70\cdot40=2800$~min $=46{,}7$~h, más 2~h: $48{,}7$~h.
  \item Ahorro: $266{,}7/48{,}7=5{,}5$ veces.
  \item Riesgos: el sustituto puede ser inexacto cerca del óptimo (se mitiga con validación, \cref{tema:6}, y con las 10 simulaciones de confirmación); puede faltar resolución si el óptimo es estrecho; y el análisis vale solo dentro del dominio muestreado.
\end{enumerate}
\resultado \textbf{48,7~h frente a 266,7~h; compensa si se valida el sustituto y se confirman los resultados.}
\comprobacion el criterio coincide con el de Martins y Ning~\cite[§10.1]{mn}: el sustituto compensa cuando su coste es pequeño frente al de las evaluaciones ahorradas.
\end{solucion}

\begin{solucion}{eje:1-jensen}
\datos $y_i=f(x_i)+\varepsilon_i$, $\E\varepsilon_i=0$.
\planteamiento la función $m(v)=\min_iv_i$ es cóncava; Jensen para cóncavas: $\E[m(V)]\le m(\E V)$.
\pasos
\begin{enumerate}[label=(\alph*)]
  \item $m$ es cóncava porque es el mínimo de funciones lineales $v\mapsto v_i$. Por Jensen, $\E[\min_iy_i]\le\min_i\E[y_i]=\min_if(x_i)$.
  \item Sean $M_n=\max_{i\le n}Z_i$. Para cada realización, $M_{n+1}=\max(M_n,Z_{n+1})\ge M_n$; tomando esperanzas, $e_{n+1}\ge e_n$.
  \item Con $n$ ajustes iguales, el sesgo de «mejor observado» es $\sigma e_n$, que no disminuye al aumentar $n$: más ensayos sin modelo del ruido no corrigen, sino que agravan, el optimismo.
\end{enumerate}
\resultado \textbf{$\E[\min_iy_i]\le\min_if(x_i)$; $e_n$ es no decreciente; el sesgo crece con el número de candidatos.}
\comprobacion con los valores del texto, $e_2=0{,}564<e_5=1{,}163<e_{10}=1{,}539<e_{30}=2{,}043$.
\end{solucion}

\begin{solucion}{eje:1-horno}
\datos horno con tres parámetros; un día por ensayo; dureza $45\pm2$~HRC; maximizar tenacidad.
\planteamiento recorrer la \cref{def:restricciones}.
\pasos
\begin{enumerate}[label=\arabic*.]
  \item Variables: $T_a\in[840,900]$~°C, $t\in[20,60]$~min (posiblemente logarítmica), $T_r\in[400,600]$~°C.
  \item Objetivo: minimizar $f=-\text{energía Charpy media}$ (J).
  \item Restricciones: $c_1=43-\text{HRC}\le0$ y $c_2=\text{HRC}-47\le0$, caras (se miden en el ensayo) y del mismo experimento que el objetivo; rangos de temperatura, conocidos; seguridad: no superar la temperatura máxima del horno.
  \item Ruido: dispersión entre probetas (la Charpy suele dispersar bastante) y de medida; se promedian las 5 probetas.
  \item Medida: regret simple (se busca la receta final).
  \item Presupuesto razonable: 20–30 ensayos (un mes), con un diseño inicial de unos 8.
  \item Método: tres variables y decenas de ensayos ruidosos: optimización bayesiana con restricciones, precedida de un factorial $2^3$ con punto central para cribar.
\end{enumerate}
\resultado \textbf{Problema restringido de 3 variables, ruido moderado, regret simple; optimización bayesiana con EI restringida tras un cribado factorial.}
\comprobacion es la misma estructura que la inyectora, con la restricción en forma de banda (dos desigualdades).
\end{solucion}

% ---------------------------------------------------------------------
\section{Tema 2}

\begin{solucion}{eje:2-matriz}
\datos $2^3$ en orden estándar.
\planteamiento columnas $A$, $B$, $C$ según la \cref{def:factorial}; interacciones por productos.
\pasos
\begin{enumerate}[label=\arabic*.]
  \item Filas $(A,B,C)$: $(-,-,-)$, $(+,-,-)$, $(-,+,-)$, $(+,+,-)$, $(-,-,+)$, $(+,-,+)$, $(-,+,+)$, $(+,+,+)$.
  \item $AC=(+,-,+,-,-,+,-,+)$; $ABC=(-,+,+,-,+,-,-,+)$.
  \item $AC\cdot A=\sum_iu_{iA}^2u_{iC}=\sum_iu_{iC}=0$; $AC\cdot ABC=\sum_iu_{iB}=0$.
  \item El orden estándar hace que $C$ cambie una sola vez: cualquier deriva lenta (calentamiento) se confundiría con el efecto de $C$. Aleatorizar reparte la deriva entre todos los efectos como ruido.
\end{enumerate}
\resultado \textbf{Matriz anterior; las columnas son ortogonales; se aleatoriza para no confundir derivas con efectos.}
\comprobacion las columnas son las del factorial de la práctica (Parte~2) y cumplen $U^\top U=8I$.
\end{solucion}

\begin{solucion}{eje:2-plasma}
\datos $\hat E_A=-101{,}625$, $\hat E_C=306{,}125$, $\hat E_{AC}=-153{,}625$; $n=2$, $N=8$; $MS_E=1738$ con 12 gl; media $776{,}06$.
\planteamiento \cref{prop:anova}: $SS=\frac{nN}{4}\hat E^2=4\hat E^2$; coeficientes $=\hat E/2$ (\cref{cor:coeficientes}).
\pasos
\begin{enumerate}[label=\arabic*.]
  \item $SS_A=4\cdot101{,}625^2=4\cdot10\,327{,}64=41\,310{,}6$.
  \item $SS_C=4\cdot306{,}125^2=4\cdot93\,712{,}5=374\,850{,}1$.
  \item $SS_{AC}=4\cdot153{,}625^2=4\cdot23\,600{,}6=94\,402{,}6$.
  \item $F_{AC}=94\,402{,}6/1738=54{,}32>4{,}75$: significativa.
  \item $\hat y=776{,}06-50{,}81\,u_A+153{,}06\,u_C-76{,}81\,u_Au_C$.
\end{enumerate}
\resultado \textbf{$SS_A=41\,311$, $SS_C=374\,850$, $SS_{AC}=94\,403$; $F_{AC}=54{,}3$, significativa; modelo anterior.}
\comprobacion $SS_A+SS_C=416\,161$ y $SS_{AC}=94\,403$ coinciden con la tabla ANOVA del modelo reducido de STAT~503~\cite{psu503}, y $F_{AC}=54{,}31$ también.
\end{solucion}

\begin{solucion}{eje:2-alias}
\datos generadores $D=AB$, $E=AC$.
\planteamiento \cref{def:fraccion}: palabras $I=ABD$, $I=ACE$ y su producto.
\pasos
\begin{enumerate}[label=\arabic*.]
  \item $ABD\cdot ACE=A^2BCDE=BCDE$. Relación definidora: $I=ABD=ACE=BCDE$.
  \item La palabra más corta tiene 3 letras: resolución III.
  \item Alias de $A$: $A\cdot ABD=BD$, $A\cdot ACE=CE$, $A\cdot BCDE=ABCDE$: $A=BD=CE=ABCDE$.
  \item Alias de $BC$: $BC\cdot ABD=ACD$, $BC\cdot ACE=ABE$, $BC\cdot BCDE=DE$: $BC=DE=ACD=ABE$.
\end{enumerate}
\resultado \textbf{$I=ABD=ACE=BCDE$; resolución III; $A=BD=CE$, $BC=DE$.}
\comprobacion con 8 ensayos y 5 factores no puede haber resolución IV sin que alguna palabra tenga 3 letras; el cribado es válido solo si se aceptan los alias de los principales con dobles, lo que en la inyectora, con la interacción $AC$ (temperatura--enfriamiento) conocida, sería arriesgado.
\end{solucion}

\begin{solucion}{eje:2-canonico}
\datos $\hat y=80+4u_1+8u_2-3u_1^2-5u_2^2+2u_1u_2$.
\planteamiento \cref{prop:estacionario}: $b=(4,8)$, $B=\bigl(\begin{smallmatrix}-3&1\\1&-5\end{smallmatrix}\bigr)$.
\pasos
\begin{enumerate}[label=\arabic*.]
  \item $\det B=15-1=14$; $B^{-1}=\frac1{14}\bigl(\begin{smallmatrix}-5&-1\\-1&-3\end{smallmatrix}\bigr)$.
  \item $B^{-1}b=\frac1{14}(-20-8;\ -4-24)=(-2;-2)$; $u_s=-\frac12(-2;-2)=(1;1)$.
  \item $\hat y(u_s)=80+\frac12(4\cdot1+8\cdot1)=86$.
  \item Autovalores: $\lambda^2+8\lambda+14=0$, $\lambda=-4\pm\sqrt2$: $-2{,}586$ y $-5{,}414$, ambos negativos.
\end{enumerate}
\resultado \textbf{Máximo en $u_s=(1,1)$ con $\hat y=86\,\%$.}
\comprobacion $\hat y(1,1)=80+4+8-3-5+2=86$ y el gradiente $(4-6+2;\ 8-10+2)=(0;0)$.
\end{solucion}

\begin{solucion}{eje:2-ccd}
\datos CCD rotable; $k=3$ con 6 centrales; $k=4$ con 4 centrales.
\planteamiento \cref{def:ccd}: $2^k+2k+n_c$ ensayos, $\alpha=(2^k)^{1/4}$, $p=1+k+k(k+1)/2$.
\pasos
\begin{enumerate}[label=\arabic*.]
  \item $k=3$: $8+6+6=20$ ensayos; $\alpha=8^{1/4}=1{,}682$; $p=1+3+6=10$; quedan $10$ gl.
  \item $k=4$: $16+8+4=28$ ensayos; $\alpha=16^{1/4}=2$; $p=1+4+10=15$; quedan $13$ gl.
\end{enumerate}
\resultado \textbf{20 ensayos y $\alpha=1{,}682$; 28 ensayos y $\alpha=2$; 10 y 15 coeficientes; 10 y 13 gl de error.}
\comprobacion con $k=2$ se obtiene $\alpha=\sqrt2$, el valor del \cref{ej:ccd}.
\end{solucion}

\begin{solucion}{eje:2-lhs}
\datos $\pi_1=(0,1,2,3,4)$, $\pi_2=(0,2,4,1,3)$, $v=0{,}5$.
\planteamiento $z_{ij}=(\pi_j(i)+0{,}5)/5$.
\pasos
\begin{enumerate}[label=\arabic*.]
  \item Puntos: $(0{,}1;0{,}1)$, $(0{,}3;0{,}5)$, $(0{,}5;0{,}9)$, $(0{,}7;0{,}3)$, $(0{,}9;0{,}7)$.
  \item Cada coordenada toma los valores $0{,}1;0{,}3;0{,}5;0{,}7;0{,}9$, uno en cada estrato de anchura $0{,}2$: es un LHS.
  \item Distancias más cortas: entre puntos consecutivos en la primera coordenada con salto 2 o 3 en la segunda, $\sqrt{0{,}2^2+0{,}4^2}=\sqrt{0{,}2}=0{,}4472$ (por ejemplo, el primero y el segundo). Las demás son $0{,}6325$, $0{,}8944$ o $1$.
  \item LHS diagonal: puntos $(0{,}1;0{,}1),\dots,(0{,}9;0{,}9)$; distancia mínima $\sqrt{0{,}2^2+0{,}2^2}=0{,}2828$.
\end{enumerate}
\resultado \textbf{Distancia mínima $0{,}447$ frente a $0{,}283$ del diagonal.}
\comprobacion una búsqueda exhaustiva sobre las 120 permutaciones (con $\pi_1$ identidad) confirma que $0{,}447$ es el máximo alcanzable con centros de celda.
\end{solucion}

\begin{solucion}{eje:2-doptimo}
\datos (a) filas de $X=[\mathbf 1\ u_1\ u_2]$: $(1,-1,-1)$, $(1,1,-1)$, $(1,-1,1)$ y $(1;\,0{,}5;\,0{,}5)$.
\planteamiento calcular $X^\top X$, invertir y comparar con $1/N$.
\pasos
\begin{enumerate}[label=(\alph*)]
  \item $X^\top X=\begin{pmatrix}4&-0{,}5&-0{,}5\\-0{,}5&3{,}25&-0{,}75\\-0{,}5&-0{,}75&3{,}25\end{pmatrix}$; invirtiendo, $[(X^\top X)^{-1}]_{11}=0{,}3355$ (fila de $\beta_1$), mayor que $1/4$: la varianza de $\hat\beta_1$ aumenta un $34\,\%$.
  \item Si $\Var(\hat\beta_j)=\sigma^2/N$, las dos desigualdades de la demostración son igualdades. La segunda, $A_{jj}=N$, exige $u_{ij}^2=1$ para todo $i$. La primera es Cauchy--Schwarz en $\langle\cdot,\cdot\rangle_A$, con igualdad si y solo si $A^{-1}e_j$ es proporcional a $e_j$, es decir, $e_j$ es autovector de $A$: la columna $j$ de $A$ es $A_{jj}e_j$, lo que significa que la columna $j$ de $X$ es ortogonal a las demás. Recíprocamente, si se cumplen ambas, $A$ tiene la fila $j$ igual a $Ne_j^\top$ y $[A^{-1}]_{jj}=1/N$.
\end{enumerate}
\resultado \textbf{$\Var(\hat\beta_1)=0{,}3355\,\sigma^2>\sigma^2/4$; la cota se alcanza si y solo si la columna es de $\pm1$ y ortogonal a las demás.}
\comprobacion en el $2^2$ completo, $X^\top X=4I$ y $[(X^\top X)^{-1}]_{11}=0{,}25$.
\end{solucion}

\begin{solucion}{eje:2-sinter}
\datos $2^3$ sin réplicas; modelo con $1,A,B,C$.
\planteamiento por la \cref{prop:ortogonal}, $\bar y=\sum_S\hat\beta_Sm_S$ con todas las columnas, y los $\hat\beta$ no cambian al quitar términos (\cref{cor:coeficientes}).
\pasos
\begin{enumerate}[label=\arabic*.]
  \item Con todos los términos, el ajuste es exacto (8 coeficientes, 8 datos): $y=\hat\beta_0\mathbf 1+\sum_{S}\hat\beta_Sm_S$.
  \item El modelo reducido predice $\hat y=\hat\beta_0\mathbf 1+\hat\beta_Am_A+\hat\beta_Bm_B+\hat\beta_Cm_C$; el residuo es $\sum_{S\in\{AB,AC,BC,ABC\}}\hat\beta_Sm_S$.
  \item Por ortogonalidad, $\norm{\text{residuo}}^2=\sum_S\hat\beta_S^2\cdot8=\sum_S SS_S$ (con $n=1$, $SS_S=2\hat E_S^2=8\hat\beta_S^2$).
  \item Si $AC$ es grande, su $SS$ entra en el «error»: $MS_E$ se infla, los $F$ de los efectos principales bajan y puede declararse no significativo un efecto real. En el \cref{ej:factorial}, $SS_{AC}=0{,}0039$ es 15 veces el error puro de las réplicas.
\end{enumerate}
\resultado \textbf{$SS_{\text{res}}=SS_{AB}+SS_{AC}+SS_{BC}+SS_{ABC}$; una interacción grande en el residuo resta potencia al contraste.}
\comprobacion los grados de libertad cuadran: $7=3$ (principales) $+4$ (residuo).
\end{solucion}

% ---------------------------------------------------------------------
\section{Tema 3}

\begin{solucion}{eje:3-nucleos}
\planteamiento elegir núcleos según la \cref{def:nucleos} y las reglas de combinación.
\pasos
\begin{enumerate}[label=(\alph*)]
  \item Matérn 5/2 (o exponencial cuadrático) con una longitud de escala grande frente al rango de $a/c$, más una media lineal o un núcleo lineal sumado que recoja la tendencia; $\sigma_f$ es la amplitud de las desviaciones respecto a la tendencia.
  \item Suma de un núcleo periódico con periodo 24~h (forma del ciclo diario) y un exponencial cuadrático de escala larga (tendencia); el producto de ambos permitiría que el ciclo cambiara lentamente.
  \item Matérn 5/2 con ARD: tres longitudes de escala, y se espera $\ell_P$ grande. Comparar la verosimilitud marginal con y sin $P$ confirma la sospecha.
\end{enumerate}
\resultado \textbf{(a) Matérn de escala larga con tendencia; (b) periódico más exponencial cuadrático; (c) Matérn con ARD.}
\comprobacion en la práctica, el núcleo ajustado del coste tiene longitudes de escala $(0{,}48;\ 0{,}56;\ 0{,}34)$ en una repetición, del mismo orden: con 20 puntos no se distingue que $P$ influya menos, lo que muestra los límites de ARD con pocos datos.
\end{solucion}

\begin{solucion}{eje:3-gp1}
\datos $k(x,x')=e^{-(x-x')^2/2}$; datos $(0,1)$ y después $(1,0)$; $\sigma_n\in\{0;0{,}1\}$.
\planteamiento \cref{teo:posterior}; con un dato, $\mu=k(x,0)\,y/(1+\sigma_n^2)$ y $\sigma^2=1-k(x,0)^2/(1+\sigma_n^2)$.
\pasos
\begin{enumerate}[label=\arabic*.]
  \item Sin ruido: en $x=0$, $\mu=1$, $\sigma^2=0$; en $x=1$, $k=e^{-1/2}=0{,}6065$: $\mu=0{,}6065$, $\sigma^2=1-0{,}3679=0{,}6321$.
  \item Con $\sigma_n=0{,}1$: en $x=0$, $\mu=1/1{,}01=0{,}9901$, $\sigma^2=1-1/1{,}01=0{,}0099$; en $x=1$, $\mu=0{,}6065/1{,}01=0{,}6005$, $\sigma^2=1-0{,}3679/1{,}01=0{,}6358$.
  \item Dos datos con ruido: $\Ky=\bigl(\begin{smallmatrix}1{,}01&0{,}6065\\0{,}6065&1{,}01\end{smallmatrix}\bigr)$, $\det=0{,}6522$, $\alpha=\Ky^{-1}y=(1{,}5486;-0{,}9299)$.
  \item En $x=0{,}5$: $k_*=(0{,}8825;0{,}8825)$, $\mu=0{,}8825\cdot0{,}6186=0{,}5459$; $\sigma^2=0{,}0365$.
  \item En $x=2$: $k_*=(e^{-2};e^{-1/2})=(0{,}1353;0{,}6065)$, $\mu=0{,}2096-0{,}5640=-0{,}3545$; $\sigma^2=0{,}5546$.
\end{enumerate}
\resultado \textbf{Con un dato, $(\mu,\sigma^2)$ vale $(1;0)$ y $(0{,}6065;0{,}6321)$ sin ruido, $(0{,}9901;0{,}0099)$ y $(0{,}6005;0{,}6358)$ con ruido. Con dos datos y ruido: $\mu(0{,}5)=0{,}5459$, $\sigma^2=0{,}0365$; $\mu(2)=-0{,}3545$, $\sigma^2=0{,}5546$.}
\comprobacion sin ruido el \cref{ej:gp2} da $\mu(0{,}5)=0{,}5493$ y $\sigma^2=0{,}0305$; el ruido acerca la media a cero y aumenta un poco la varianza, como cabe esperar.
\end{solucion}

\begin{solucion}{eje:3-press}
\datos $(0,1)$, $(1,2)$, $(2,2)$, $(3,5)$.
\planteamiento mínimos cuadrados, $h_{ii}=\frac1n+\frac{(x_i-\bar x)^2}{\sum(x_l-\bar x)^2}$ para la recta, y \cref{prop:loo-lineal}.
\pasos
\begin{enumerate}[label=\arabic*.]
  \item $\bar x=1{,}5$, $\bar y=2{,}5$, $S_{xx}=5$, $S_{xy}=(-1{,}5)(-1{,}5)+(-0{,}5)(-0{,}5)+0{,}5\cdot(-0{,}5)+1{,}5\cdot2{,}5=6$; pendiente $1{,}2$, ordenada $2{,}5-1{,}8=0{,}7$.
  \item Residuos: $1-0{,}7=0{,}3$; $2-1{,}9=0{,}1$; $2-3{,}1=-1{,}1$; $5-4{,}3=0{,}7$.
  \item $h_{ii}=0{,}25+(x_i-1{,}5)^2/5$: $0{,}7;\ 0{,}3;\ 0{,}3;\ 0{,}7$.
  \item $e_{(i)}=e_i/(1-h_{ii})$: $1{,}0;\ 0{,}1429;\ -1{,}5714;\ 2{,}3333$.
  \item PRESS medio: $(1+0{,}0204+2{,}4694+5{,}4444)/4=2{,}2336$; raíz $1{,}495$.
  \item Comprobación directa de $i=4$: sin $(3,5)$, la recta por $(0,1),(1,2),(2,2)$ tiene pendiente $0{,}5$ y ordenada $1{,}1667$; predice $2{,}6667$ en $x=3$ y el residuo es $2{,}3333$.
\end{enumerate}
\resultado \textbf{Recta $\hat y=0{,}7+1{,}2x$; residuos LOO $(1{,}0;\ 0{,}143;\ -1{,}571;\ 2{,}333)$; raíz del PRESS medio $1{,}495$.}
\comprobacion la traza de $H$ es $0{,}7+0{,}3+0{,}3+0{,}7=2$, el número de coeficientes.
\end{solucion}

\begin{solucion}{eje:3-lml}
\datos $x=(0,1)$, $y=(1;0{,}9)$, $\sigma_f=1$, $\sigma_n=0{,}1$.
\planteamiento evaluar los tres términos de~\eqref{eq:lml}; para $2\times2$, $y^\top\Ky^{-1}y=(a y_2^2-2b y_1y_2+a y_1^2)/(a^2-b^2)$ con $a=1{,}01$, $b=e^{-1/(2\ell^2)}$.
\pasos
\begin{enumerate}[label=\arabic*.]
  \item $\ell=0{,}1$: $b\approx0$; $y^\top\Ky^{-1}y=1{,}81/1{,}01=1{,}7921$; $\det=1{,}0201$. Términos: $-0{,}8960$, $-0{,}0100$, $-1{,}8379$; total $-2{,}7439$.
  \item $\ell=1$: $b=0{,}6065$; $\det=0{,}6522$; $y^\top\Ky^{-1}y=1{,}1290$. Términos: $-0{,}5645$, $+0{,}2137$, $-1{,}8379$; total $-2{,}1887$.
  \item $\ell=3$: $b=0{,}9460$; $\det=0{,}1253$; $y^\top\Ky^{-1}y=1{,}0009$. Términos: $-0{,}5004$, $+1{,}0387$, $-1{,}8379$; total $-1{,}2996$.
\end{enumerate}
\resultado \textbf{Se elige $\ell=3$, con $\log p=-1{,}300$.}
\comprobacion los datos casi coinciden ($1$ y $0{,}9$): una función muy suave los explica bien (ajuste $-0{,}50$) y además concentra la probabilidad en pocas funciones (término de complejidad positivo y grande), mientras que $\ell=0{,}1$ trata los dos datos como independientes y paga en ajuste. Es la navaja de Occam comentada tras la \cref{prop:lml}.
\end{solucion}

\begin{solucion}{eje:3-jitter}
\datos $K=\bigl(\begin{smallmatrix}1&\rho\\\rho&1\end{smallmatrix}\bigr)$ con $\rho=e^{-\delta^2/2}$.
\planteamiento autovalores de una matriz simétrica $2\times2$ con diagonal constante.
\pasos
\begin{enumerate}[label=\arabic*.]
  \item $\det(K-\lambda I)=(1-\lambda)^2-\rho^2=0\Rightarrow\lambda=1\pm\rho$, con autovectores $(1,\pm1)/\sqrt2$.
  \item $\delta=0{,}01$: $1-\rho=5{,}0\cdot10^{-5}$; número de condición $(1+\rho)/(1-\rho)=4{,}0\cdot10^4$.
  \item $\delta=0{,}001$: $1-\rho=5{,}0\cdot10^{-7}$; número de condición $4{,}0\cdot10^6$.
  \item Con muchos puntos próximos, el autovalor mínimo se acerca al error de redondeo de doble precisión (módulo~11, tema~5) y la factorización de Cholesky falla o da raíces de números negativos; sumar $10^{-6}I$ eleva todos los autovalores en $10^{-6}$ y acota el número de condición en torno a $2\cdot10^6$.
\end{enumerate}
\resultado \textbf{Autovalores $1\pm e^{-\delta^2/2}$; condición $4\cdot10^4$ y $4\cdot10^6$; el \emph{jitter} evita la casi singularidad.}
\comprobacion para $\delta\to0$, $1-\rho\approx\delta^2/2$: $0{,}01^2/2=5\cdot10^{-5}$.
\end{solucion}

\begin{solucion}{eje:3-mezcla}
\datos $\mu_m=(1;2;3)$, $\sigma_m=0{,}5$.
\planteamiento \cref{prop:mezcla}.
\pasos
\begin{enumerate}[label=\arabic*.]
  \item $\mu=(1+2+3)/3=2$.
  \item Parte aleatoria: $\frac13(0{,}25\cdot3)=0{,}25$.
  \item Parte epistémica: $\frac13(1+0+1)=0{,}6667$.
  \item $\sigma^2=0{,}9167$, $\sigma=0{,}9574$.
  \item Si las redes coincidieran, la parte epistémica sería nula y $\sigma=0{,}5$.
\end{enumerate}
\resultado \textbf{$\mu=2$, $\sigma=0{,}957$ (aleatoria $0{,}25$, epistémica $0{,}667$); con redes coincidentes, $\sigma=0{,}5$.}
\comprobacion con la primera forma: $\frac13\bigl((0{,}25+1)+(0{,}25+4)+(0{,}25+9)\bigr)-4=4{,}9167-4=0{,}9167$.
\end{solucion}

\begin{solucion}{eje:3-monotona}
\datos $f\sim\GP(0,k)$, datos $\D_n$ y una observación nueva $y_{n+1}$.
\planteamiento \cref{lem:condicion}: la varianza condicional es la previa menos una forma cuadrática semidefinida positiva.
\pasos
\begin{enumerate}[label=(\alph*)]
  \item $\sigma_n^2(x)=k(x,x)-k_*^\top\Ky^{-1}k_*$, y $\Ky^{-1}$ es definida positiva, luego $k_*^\top\Ky^{-1}k_*\ge0$ y $\sigma_n^2(x)\le k(x,x)$.
  \item Dados $\D_n$, el par $(f(x),y_{n+1})$ es normal con varianzas $\sigma_n^2(x)$, $s^2=\sigma_n^2(x_{n+1})+\sigma_\varepsilon^2>0$ y covarianza posterior $c_n(x,x_{n+1})$. Condicionando a $y_{n+1}$ con el lema: $\sigma_{n+1}^2(x)=\sigma_n^2(x)-c_n(x,x_{n+1})^2/s^2\le\sigma_n^2(x)$.
\end{enumerate}
\resultado \textbf{La varianza posterior no supera la previa y no aumenta al añadir observaciones.}
\comprobacion en el \cref{eje:4-creyente}, $\sigma^2(0{,}5)$ baja de $0{,}2212$ a $0{,}0305$ al añadir el punto $x=1$.
\end{solucion}

\begin{solucion}{eje:3-optimos}
\datos dos máximos de la verosimilitud marginal con $\sigma_n=0{,}1$: $\ell=0{,}96$ (local) y $\ell=0{,}03$ (global).
\planteamiento interpretar cada régimen con la \cref{obs:cholesky}.
\pasos
\begin{enumerate}[label=(\alph*)]
  \item Con $\ell=0{,}03$, los datos están a más de seis longitudes de escala unos de otros y son casi independientes: la media pasa por cada dato y vuelve a la media previa entre ellos; las bandas son estrechas en los datos y anchísimas entre ellos. Con $\ell=0{,}96$, la media es una curva suave que no pasa exactamente por los datos y las bandas son moderadas en todo el rango.
  \item (1) Validación cruzada dejando uno fuera (\cref{teo:loo}): el modelo con $\ell=0{,}03$ predice muy mal cada dato a partir de los demás. (2) Conocimiento físico: el alabeo no oscila en décimas de segundo de enfriamiento; también puede ponerse una distribución previa sobre $\ell$ que penalice valores menores que la resolución de la máquina.
  \item Si $\sigma_n$ se estima, la diferencia entre los lotes casi coincidentes se explica como ruido, con un coste de complejidad menor que el de una $\ell$ diminuta; con el valor correcto ($0{,}165$) el máximo es único en $\ell=0{,}87$.
\end{enumerate}
\resultado \textbf{$\ell=0{,}03$ interpola el ruido; $\ell=0{,}96$ sigue la tendencia; se elige con LOO o con conocimiento previo, y el problema desaparece estimando $\sigma_n$.}
\comprobacion la práctica (Parte~3) muestra ambas curvas de $\log p$ frente a $\ell$ y confirma los valores.
\end{solucion}

\begin{solucion}{eje:3-pinn}
\datos $T(t)=40+190e^{-0{,}1t}$, $k=0{,}12$, $T_m=40$.
\planteamiento $\mathcal R=\dot T+k(T-T_m)$.
\pasos
\begin{enumerate}[label=\arabic*.]
  \item $e^{-0{,}5}=0{,}6065$; $T(5)=40+115{,}24=155{,}24$~°C; $\dot T(5)=-19\cdot0{,}6065=-11{,}52$~°C/s.
  \item $\mathcal R(5)=-11{,}52+0{,}12\cdot115{,}24=-11{,}52+13{,}83=2{,}31$~°C/s.
  \item En general $\mathcal R=(k-0{,}1)\cdot190e^{-0{,}1t}$, nulo para todo $t$ si y solo si $k=0{,}1$.
  \item Tiempo hasta 90~°C: $\ln(190/50)=1{,}335$; con $k=0{,}1$, $13{,}35$~s; con $k=0{,}12$, $11{,}13$~s.
\end{enumerate}
\resultado \textbf{$\mathcal R(5)=2{,}31$; la función es solución exacta con $k=0{,}1$; 13,35~s frente a 11,13~s.}
\comprobacion $(0{,}12-0{,}1)\cdot115{,}24=2{,}30$, igual que el paso 2 salvo redondeo.
\end{solucion}

% ---------------------------------------------------------------------
\section{Tema 4}

\begin{solucion}{eje:4-pi}
\planteamiento usar las definiciones de PI, EI (\cref{teo:ei}) y LCB.
\pasos
\begin{enumerate}[label=(\alph*)]
  \item Junto al mejor punto, $\mu\approx\fmejor$ y $\sigma$ pequeña dan $\PI\approx\Phi(0^+)\approx0{,}5$ o más si $\mu$ es algo menor; lejos, $\mu$ suele ser mayor y la PI pequeña aunque $\sigma$ sea grande. La PI recompensa la probabilidad, no la magnitud de la mejora.
  \item Sin ruido, en un punto evaluado $\sigma=0$ y $\mu=y_i\ge\fmejor$: $\EI=\max(\fmejor-y_i,0)=0$. Con la varianza predictiva, $\sigma\ge\sigma_\varepsilon>0$ y $\EI\ge\sigma_\varepsilon\varphi(z)>0$: el punto puede volver a elegirse.
  \item $\sqrt{\beta_n}$ es el número de desviaciones típicas que se resta: el nivel de optimismo. Con $\beta_n\to0$ se minimiza la media (explotación pura); con $\beta_n\to\infty$, se maximiza $\sigma_n$ (exploración pura).
\end{enumerate}
\resultado \textbf{La PI es miope respecto a la magnitud; la EI se anula en puntos evaluados sin ruido; $\beta_n$ regula explorar frente a explotar.}
\comprobacion en el \cref{ej:ei}, la PI elige $A$ con $0{,}977$ aunque su mejora esperada sea solo de dos céntimos.
\end{solucion}

\begin{solucion}{eje:4-calculo}
\datos $\fmejor=1$; $P(0{,}98;0{,}02)$, $Q(1{,}05;0{,}10)$, $R(1{,}20;0{,}30)$.
\planteamiento $z=(\fmejor-\mu)/\sigma$; $\PI=\Phi(z)$; $\EI=(\fmejor-\mu)\Phi(z)+\sigma\varphi(z)$.
\pasos
\begin{enumerate}[label=\arabic*.]
  \item $P$: $z=1$; $\PI=0{,}8413$; $\EI=0{,}02\cdot0{,}8413+0{,}02\cdot0{,}2420=0{,}02167$; $\LCB=0{,}94$.
  \item $Q$: $z=-0{,}5$; $\PI=0{,}3085$; $\EI=-0{,}05\cdot0{,}3085+0{,}10\cdot0{,}3521=0{,}01978$; $\LCB=0{,}85$.
  \item $R$: $z=-2/3$; $\PI=0{,}2525$; $\EI=-0{,}20\cdot0{,}2525+0{,}30\cdot0{,}3194=0{,}04532$; $\LCB=0{,}60$.
\end{enumerate}
\resultado \textbf{PI elige $P$; EI y LCB eligen $R$ ($\EI=0{,}0453$, $\LCB=0{,}60$).}
\comprobacion a diferencia del \cref{ej:ei}, aquí la incertidumbre de $R$ es tan grande que la EI también la prefiere: la EI explora cuando $\sigma$ compensa la peor media (\cref{prop:ei}, iii).
\end{solucion}

\begin{solucion}{eje:4-restringida}
\datos $\EI(A)=0{,}02008$, $\EI(B)=0{,}01187$; $A$: $\mu_w=0{,}44$; $B$: $\mu_w=0{,}40$; $\sigma_w=0{,}02$; límite $0{,}45$.
\planteamiento \cref{prop:eic}: $\EI_c=\EI\cdot\Phi((0{,}45-\mu_w)/\sigma_w)$.
\pasos
\begin{enumerate}[label=\arabic*.]
  \item $A$: $\Phi(0{,}5)=0{,}6915$; $\EI_c=0{,}02008\cdot0{,}6915=0{,}01389$.
  \item $B$: $\Phi(2{,}5)=0{,}9938$; $\EI_c=0{,}01187\cdot0{,}9938=0{,}01179$.
\end{enumerate}
\resultado \textbf{$\EI_c(A)=0{,}0139>\EI_c(B)=0{,}0118$: la elección no cambia, aunque la ventaja de $A$ se reduce de un $69\,\%$ a un $18\,\%$.}
\comprobacion $0{,}013\,888$ coincide con la comprobación del ejercicio~11 de la práctica.
\end{solucion}

\begin{solucion}{eje:4-creyente}
\datos núcleo exponencial cuadrático, $\sigma_f=\ell=1$; dato en $0$; dato ficticio en $1$.
\planteamiento \cref{teo:posterior}; la varianza no depende de los valores (\cref{obs:posterior}).
\pasos
\begin{enumerate}[label=(\alph*)]
  \item $\sigma_1^2(x)=1-e^{-x^2}$: $\sigma_1^2(0{,}5)=1-e^{-0{,}25}=0{,}2212$; $\sigma_1^2(1)=1-e^{-1}=0{,}6321$.
  \item Con datos en $0$ y $1$: $\sigma^2(1)=0$ (punto observado sin ruido) y $\sigma^2(0{,}5)=0{,}0305$ (\cref{ej:gp2}).
  \item La varianza solo depende de las posiciones; el valor ficticio cambia la media, no la incertidumbre. Como $\sigma^2(1)=0$, la EI en $x=1$ es $\max(\fmejor-\mu(1),0)$, y con el creyente $\mu(1)$ es la media anterior, por lo que en general la EI allí es nula o pequeña y el segundo punto se aleja.
\end{enumerate}
\resultado \textbf{$\sigma^2(0{,}5)$ baja de $0{,}2212$ a $0{,}0305$ y $\sigma^2(1)$ de $0{,}6321$ a $0$; el dato ficticio aparta el siguiente punto del lote.}
\comprobacion el resultado (b) coincide con el del \cref{ej:gp2}, calculado con un valor real en $x=1$: la varianza es la misma.
\end{solucion}

\begin{solucion}{eje:4-pareto}
\datos seis puntos (defectos \%, ciclo s); $r=(10;35)$.
\planteamiento \cref{def:pareto} y \cref{def:hv}; en 2D, ordenar por el primer objetivo y sumar rectángulos.
\pasos
\begin{enumerate}[label=\arabic*.]
  \item $(4;27)$ está dominado por $(3;26)$; $(2{,}5;32)$, por $(2;30)$. No dominados: $(2;30)$, $(3;26)$, $(5;22)$, $(8;21)$.
  \item Rectángulos: $(10-2)(35-30)=40$; $(10-3)(30-26)=28$; $(10-5)(26-22)=20$; $(10-8)(22-21)=2$. Total: 90.
  \item $(4;23)$ no está dominado. Nuevo orden: $(2;30)$, $(3;26)$, $(4;23)$, $(5;22)$, $(8;21)$: $40+28+(10-4)(26-23)+(10-5)(23-22)+2=40+28+18+5+2=93$.
\end{enumerate}
\resultado \textbf{Cuatro puntos no dominados; $\HV=90$; el nuevo punto lo eleva a 93 (+3).}
\comprobacion el aumento es el rectángulo $[4,5]\times[23,26]$ que solo domina el punto nuevo: $1\cdot3=3$. La práctica comprueba $\HV=90$.
\end{solucion}

\begin{solucion}{eje:4-coste}
\datos simulación: 20~min, EI $0{,}002$~€; lote: 120~min, EI $0{,}009$~€.
\planteamiento EI por unidad de coste~\cite{snoek}.
\pasos
\begin{enumerate}[label=\arabic*.]
  \item Simulación: $0{,}002/20=1{,}0\cdot10^{-4}$~€/min. Lote: $0{,}009/120=7{,}5\cdot10^{-5}$~€/min. Gana la simulación.
  \item Equivalencia: $\EI_{\text{sim}}/20=0{,}009/120\Rightarrow\EI_{\text{sim}}=0{,}0015$~€.
  \item Si la simulación está sesgada, su EI sobre el coste \emph{real} está mal calculada: el criterio puede preferir una y otra vez una fuente barata que no informa sobre la planta. Hay que modelar la relación entre fidelidades (y su sesgo) y reservar ensayos reales.
\end{enumerate}
\resultado \textbf{La simulación rinde $1{,}0\cdot10^{-4}$ frente a $7{,}5\cdot10^{-5}$~€/min; empate con EI de $0{,}0015$~€; riesgo de sesgo.}
\comprobacion $0{,}0015\cdot6=0{,}009$: seis simulaciones cuestan lo que un lote.
\end{solucion}

\begin{solucion}{eje:4-ei}
\datos $\EI=(\fmejor-\mu)\Phi(z)+\sigma\varphi(z)$, $z=(\fmejor-\mu)/\sigma$.
\planteamiento límites y derivadas usando la \cref{prop:ei}.
\pasos
\begin{enumerate}[label=(\alph*)]
  \item Si $\mu<\fmejor$, $z\to+\infty$, $\Phi(z)\to1$ y $\sigma\varphi(z)\to0$: $\EI\to\fmejor-\mu$. Si $\mu>\fmejor$, $z\to-\infty$ y ambos términos tienden a 0 (el primero porque $\Phi(z)\to0$). Si $\mu=\fmejor$, $\EI=\sigma\varphi(0)\to0$. En todos los casos, $\max(\fmejor-\mu,0)$.
  \item $\partial\EI/\partial\mu=-\Phi(z)$ y $\partial^2\EI/\partial\mu^2=\varphi(z)/\sigma>0$: convexa.
  \item En maximización, con $\Delta=\mu-\fmejor$, el cálculo del \cref{teo:ei} con $f=\mu+\sigma Z$ y la región $Z>-\Delta/\sigma$ da $\EI=\Delta\,\Phi(\Delta/\sigma)+\sigma\varphi(\Delta/\sigma)$. Si $\Delta\ge0$: $\Delta\Phi(\Delta/\sigma)=\Delta-\Delta\Phi(-\Delta/\sigma)=[\Delta]^+-\abs{\Delta}\Phi(-\abs\Delta/\sigma)$. Si $\Delta<0$: $\Delta\Phi(\Delta/\sigma)=-\abs{\Delta}\Phi(-\abs{\Delta}/\sigma)$ y $[\Delta]^+=0$. En ambos casos $\EI=[\Delta]^++\sigma\varphi(\Delta/\sigma)-\abs\Delta\,\Phi(-\abs\Delta/\sigma)$, usando que $\varphi$ es par.
\end{enumerate}
\resultado \textbf{$\EI\to\max(\fmejor-\mu,0)$; $\EI$ es convexa en $\mu$; la forma de maximización es la indicada.}
\comprobacion con $\Delta=0{,}02$, $\sigma=0{,}02$: $0{,}02\Phi(1)+0{,}02\varphi(1)=0{,}02167$, el mismo valor que $P$ en el \cref{eje:4-calculo}, como debe ser por simetría.
\end{solucion}

\begin{solucion}{eje:4-chebyshev}
\datos $F^{(1)}=(0;1)$, $F^{(2)}=(1;0)$, $F^{(3)}=(0{,}6;0{,}6)$; $z^\star=(0;0)$.
\planteamiento comprobar dominancia; comparar sumas ponderadas; buscar pesos de Chebyshev.
\pasos
\begin{enumerate}[label=(\alph*)]
  \item Ningún punto es menor o igual que otro en ambas coordenadas con alguna estricta: $F^{(3)}$ mejora a $F^{(1)}$ en la segunda y empeora en la primera, y lo análogo con $F^{(2)}$. Los tres son óptimos de Pareto.
  \item Normalizando $w_1+w_2=1$: las sumas valen $w_2$, $w_1$ y $0{,}6$. Para que $F^{(3)}$ fuera minimizador haría falta $0{,}6\le w_2$ y $0{,}6\le w_1$, es decir, $w_1+w_2\ge1{,}2$, imposible.
  \item Con $w=(1;1)$: $s(F^{(1)})=\max(0,1)=1$, $s(F^{(2)})=1$, $s(F^{(3)})=0{,}6$. $F^{(3)}$ es el único minimizador.
\end{enumerate}
\resultado \textbf{Los tres son de Pareto; ninguna suma ponderada selecciona $F^{(3)}$; Chebyshev con $w=(1,1)$ sí.}
\comprobacion $F^{(3)}$ está por encima del segmento que une $F^{(1)}$ y $F^{(2)}$ ($0{,}6+0{,}6>1$): es un punto de una parte no convexa del frente.
\end{solucion}

% ---------------------------------------------------------------------
\section{Tema 5}

\begin{solucion}{eje:5-lista}
\planteamiento aplicar los puntos 1–5 del \cref{alg:lista}.
\pasos
\begin{enumerate}[label=\arabic*.]
  \item Objetivo: maximizar el título (g/L) a 72~h; si importa el coste, título por euro de sustrato.
  \item Variables: temperatura y pH lineales; glucosa e inóculo, quizá logarítmicas (abarcan un factor 4 y 10).
  \item Restricciones: rangos conocidos (viabilidad de la cepa); posible restricción oculta (tandas que se contaminan o no crecen: fallos sin dato); seguridad, no superar la temperatura de muerte celular.
  \item Ruido: variabilidad biológica entre tandas, alta; conviene replicar el punto central para estimarla.
  \item Presupuesto: lotes de 4 (biorreactores) cada 4 días; la medida a 24~h es una fidelidad baja.
\end{enumerate}
\resultado \textbf{Problema de 4 variables, ruido biológico alto, lotes de 4 y una fidelidad baja: optimización bayesiana por lotes con restricción oculta.}
\comprobacion es el mismo esquema que el hormigón: lotes semanales y una medida temprana correlacionada.
\end{solucion}

\begin{solucion}{eje:5-fck}
\datos C25; $\sigma_p\in\{3{,}5;\ 5\}$; $f_{28}=105\cdot8^{-a/c}$.
\planteamiento \cref{def:fck}: $f_{cm}=f_{ck}+1{,}645\sigma_p$; $a/c=\ln(105/f)/\ln8$.
\pasos
\begin{enumerate}[label=\arabic*.]
  \item $\sigma_p=3{,}5$: $25+5{,}758=30{,}76$~MPa. $\sigma_p=5$: $25+8{,}225=33{,}23$~MPa.
  \item Margen ganado: $33{,}23-30{,}76=2{,}47$~MPa.
  \item $a/c$ máxima: con $33{,}23$~MPa, $\ln(3{,}160)/\ln8=1{,}1507/2{,}0794=0{,}553$; con $30{,}76$~MPa, $\ln(3{,}414)/\ln 8=1{,}2278/2{,}0794=0{,}590$.
\end{enumerate}
\resultado \textbf{30,76 y 33,23~MPa; se ganan 2,47~MPa; la $a/c$ puede subir de 0,553 a 0,590.}
\comprobacion $105\cdot8^{-0{,}590}=105/3{,}414=30{,}76$.
\end{solucion}

\begin{solucion}{eje:5-abrams}
\datos $(0{,}40;45{,}7)$, $(0{,}50;37{,}1)$, $(0{,}60;30{,}2)$.
\planteamiento $\ln f=\ln A-(a/c)\ln B$: regresión lineal de $\ln f$ sobre $a/c$ (módulo~0.9, §6.1 y §6.4).
\pasos
\begin{enumerate}[label=\arabic*.]
  \item $\ln f=(3{,}8221;\ 3{,}6136;\ 3{,}4078)$; medias $\overline{a/c}=0{,}5$ y $\overline{\ln f}=3{,}6145$.
  \item $S_{xy}=(-0{,}1)(0{,}2076)+0+(0{,}1)(-0{,}2067)=-0{,}04143$; $S_{xx}=0{,}02$; pendiente $-2{,}0713$.
  \item $\ln A=3{,}6145+2{,}0713\cdot0{,}5=4{,}6502$: $A=104{,}6$~MPa; $\ln B=2{,}0713$: $B=7{,}935$.
  \item $a/c_{\max}=(4{,}6502-\ln36{,}58)/2{,}0713=(4{,}6502-3{,}5995)/2{,}0713=0{,}507$.
  \item La ley es multiplicativa: el logaritmo la vuelve lineal y hace el error relativo (no el absoluto) homogéneo, lo que suele ajustarse mejor a la dispersión de las resistencias.
\end{enumerate}
\resultado \textbf{$f=104{,}6/7{,}935^{a/c}$; $a/c\le0{,}507$.}
\comprobacion los valores verdaderos del simulador son $A=105$ y $B=8$; el ajuste los recupera con los datos redondeados.
\end{solucion}

\begin{solucion}{eje:5-adquisicion}
\datos $g^*=335$; candidatos $(325;0{,}70)$, $(300;0{,}10)$, $(340;0{,}99)$; mezcla $(0{,}45;0{,}15;360)$.
\planteamiento \cref{obs:objconocido}: $\alpha=\max(g^*-g,0)\cdot p$.
\pasos
\begin{enumerate}[label=\arabic*.]
  \item $X_1$: $10\cdot0{,}70=7{,}0$. $X_2$: $35\cdot0{,}10=3{,}5$. $X_3$: $0\cdot0{,}99=0$.
  \item Orden: $X_1>X_2>X_3$.
  \item Mezcla: cemento $360\cdot0{,}85=306$~kg; cenizas $54$~kg. $\gwp=0{,}93\cdot306+0{,}03\cdot54+45=284{,}6+1{,}6+45=331{,}2$.
  \item Agua: $0{,}45\cdot360=162$~kg/m$^3$ (dentro de $[160,195]$).
  \item Coste: $0{,}13\cdot306+0{,}05\cdot54+55=39{,}78+2{,}70+55=97{,}48$~€/m$^3$.
\end{enumerate}
\resultado \textbf{$\alpha=(7{,}0;\ 3{,}5;\ 0)$, se elige $X_1$; la mezcla tiene $\gwp=331{,}2$, agua 162~kg/m$^3$ y coste $97{,}48$~€/m$^3$.}
\comprobacion $X_3$ tiene $\alpha=0$ aunque sea casi seguro que cumple, porque no mejora el objetivo: la adquisición no premia la factibilidad sin mejora.
\end{solucion}

\begin{solucion}{eje:5-mezclas}
\datos tres componentes con $x_1+x_2+x_3=1$, $x_i\ge0$.
\planteamiento región de mezclas~\cite[§11.6]{montgomery}.
\pasos
\begin{enumerate}[label=(\alph*)]
  \item Es el símplice de vértices $(1,0,0)$, $(0,1,0)$, $(0,0,1)$: un triángulo equilátero en coordenadas baricéntricas, con dos grados de libertad.
  \item Retícula $\{3,2\}$: los 3 vértices y los 3 puntos medios de los lados: $(\tfrac12,\tfrac12,0)$, $(\tfrac12,0,\tfrac12)$, $(0,\tfrac12,\tfrac12)$. Seis puntos.
  \item Con $x_1\ge0{,}5$ quedan $(1,0,0)$, $(\tfrac12,\tfrac12,0)$ y $(\tfrac12,0,\tfrac12)$, que son los vértices de la nueva región: un triángulo semejante con lado la mitad.
\end{enumerate}
\resultado \textbf{Símplice; seis puntos; con $x_1\ge0{,}5$ quedan tres, vértices de un subsímplice.}
\comprobacion el número de puntos de una retícula $\{q,m\}$ es $\binom{q+m-1}{m}=\binom{4}{2}=6$.
\end{solucion}

\begin{solucion}{eje:5-fidelidad}
\datos $s=0{,}20$; $f_7=22{,}0$; $\rho=f_7/f_{28}=0{,}68-0{,}45s$ con desviación $0{,}03$.
\planteamiento $f_{28}=f_7/\rho$; propagación lineal: $\sigma_{f_{28}}\approx\abs{\partial f_{28}/\partial\rho}\,\sigma_\rho=f_7\sigma_\rho/\rho^2$.
\pasos
\begin{enumerate}[label=\arabic*.]
  \item $\rho=0{,}68-0{,}09=0{,}59$; $f_{28}=22/0{,}59=37{,}29$~MPa.
  \item $\sigma_{f_{28}}=22\cdot0{,}03/0{,}3481=1{,}90$~MPa.
  \item $\Pr[f_{28}\ge36{,}58]\approx\Phi\bigl((37{,}29-36{,}58)/1{,}90\bigr)=\Phi(0{,}373)=0{,}65<0{,}9$.
\end{enumerate}
\resultado \textbf{$f_{28}\approx37{,}3\pm1{,}9$~MPa; probabilidad de cumplir $0{,}65$: no se puede decidir sin el ensayo a 28 días.}
\comprobacion sin contar siquiera el ruido de la propia medida de $f_7$ (1,5~MPa, que añadiría $1{,}5/0{,}59=2{,}5$~MPa), la incertidumbre ya es mayor que el margen.
\end{solucion}

\begin{solucion}{eje:5-ofat}
\datos (a) $f=\sum_jf_j(u_j)$; (b) $f=(u_1-u_2)^2+0{,}05(u_1+u_2-2)^2$, partida $(0,0)$.
\planteamiento (a) separabilidad; (b) minimización exacta alternada: $u_1=(1{,}9u_2+0{,}2)/2{,}1$ y $u_2=(1{,}9u_1+0{,}2)/2{,}1$.
\pasos
\begin{enumerate}[label=\arabic*.]
  \item (a) Minimizar $f$ en $u_j$ con las demás fijas es minimizar $f_j(u_j)$, cuyo minimizador no depende de las demás. Tras una pasada, cada $u_j$ minimiza su sumando y $f$ alcanza $\sum_j\min f_j$, el mínimo global.
  \item (b) Primera pasada: $u_1=0{,}2/2{,}1=0{,}0952$; $u_2=(0{,}1810+0{,}2)/2{,}1=0{,}1814$; $f=0{,}156$.
  \item Segunda pasada: $u_1=(0{,}3447+0{,}2)/2{,}1=0{,}2594$; $u_2=(0{,}4928+0{,}2)/2{,}1=0{,}3299$; $f=0{,}0050+0{,}0995=0{,}104$.
  \item (c) La hessiana es $\bigl(\begin{smallmatrix}2{,}1&-1{,}9\\-1{,}9&2{,}1\end{smallmatrix}\bigr)$, de autovalores $0{,}2$ y $4$: número de condición 20. El error $(1-u_1,1-u_2)$ se multiplica por $(1{,}9/2{,}1)^2=0{,}819$ en cada pasada: convergencia lineal lenta, como el descenso de gradiente en cuadráticas mal condicionadas (módulo~7, §2.3).
\end{enumerate}
\resultado \textbf{OFAT es exacto sin interacciones; en el valle oblicuo, tras dos pasadas está en $(0{,}259;0{,}330)$ con $f=0{,}104$.}
\comprobacion $1-u_2$ pasa de $0{,}8186$ a $0{,}6701$; el cociente es $0{,}819$.
\end{solucion}

\begin{solucion}{eje:5-fermentacion}
\datos 4 variables, 40 tandas en lotes de 4, medida a 24~h.
\planteamiento recorrer el \cref{alg:lista} con los resultados de los temas~2 a~4.
\pasos
\begin{enumerate}[label=\arabic*.]
  \item Cribado (8 tandas, 2 semanas): fracción $2^{4-1}$ de resolución IV ($D=ABC$), que estima los cuatro efectos principales libres de interacciones dobles (\cref{def:fraccion}); se añade si es posible una réplica del centro para estimar el ruido.
  \item Diseño inicial: las 8 tandas del cribado sirven ya como datos; se completan hasta 12 con 4 puntos de un LHS para cubrir el interior.
  \item Sustitutos: PG Matérn 5/2 con ARD y ruido estimado para el título a 72~h; un clasificador (o un PG sobre una variable 0/1) para el fallo de la tanda; un PG para la relación entre el título a 24 y a 72~h.
  \item Adquisición: EI ruidosa por lotes de 4 (creyente del kriging o qNEI en BoTorch) multiplicada por la probabilidad de que la tanda no falle.
  \item Baja fidelidad: a las 24~h se actualiza el modelo con la medida temprana y, si la predicción a 72~h es muy mala, se interrumpe la tanda y se libera el biorreactor (similar a la poda de Hyperband).
  \item Parada: presupuesto, o antes si la cota de la \cref{prop:cota} es menor que la desviación entre réplicas.
  \item Confirmación: 4 tandas en la recomendación (un lote completo).
  \item Reparto aproximado: 8 de cribado, 4 de relleno, 24 de optimización (6 lotes) y 4 de confirmación.
\end{enumerate}
\resultado \textbf{Plan en cuatro fases con 8 + 4 + 24 + 4 tandas, EI por lotes con restricción oculta y uso de la medida a 24~h para podar.}
\comprobacion el plan respeta el presupuesto ($8+4+24+4=40$) y el paralelismo (todas las fases en múltiplos de 4).
\end{solucion}

% ---------------------------------------------------------------------
\section{Tema 6}

\begin{solucion}{eje:6-rejilla}
\datos 9 ensayos en 2 hiperparámetros; 64 en 3.
\planteamiento la rejilla $m^d$ proyecta $m$ valores; $n$ ensayos aleatorios proyectan $n$ (con probabilidad 1).
\pasos
\begin{enumerate}[label=\arabic*.]
  \item Rejilla $3\times3$: 3 valores del hiperparámetro relevante; aleatoria: 9.
  \item 64 en 3 dimensiones: rejilla $4\times4\times4$, 4 valores; aleatoria: 64.
\end{enumerate}
\resultado \textbf{3 frente a 9 valores; 4 frente a 64.}
\comprobacion la ventaja relativa es $n^{1-1/d}$: $9^{1/2}=3$ y $64^{2/3}=16$.
\end{solucion}

\begin{solucion}{eje:6-aleatoria}
\datos (a) $q=0{,}02$, $p=0{,}95$; (b) $q=0{,}05$, $p=0{,}99$; (c) $n=27$, $q=0{,}05$.
\planteamiento \cref{prop:aleatoria}.
\pasos
\begin{enumerate}[label=\arabic*.]
  \item $\ln0{,}05/\ln0{,}98=-2{,}9957/-0{,}0202=148{,}3$: 149 ensayos.
  \item $\ln0{,}01/\ln0{,}95=-4{,}6052/-0{,}0513=89{,}8$: 90 ensayos.
  \item $1-0{,}95^{27}=1-0{,}2503=0{,}7497$.
\end{enumerate}
\resultado \textbf{149; 90; probabilidad $0{,}75$.}
\comprobacion la práctica comprueba los valores 59 y 149; $0{,}98^{149}=0{,}049<0{,}05$ y $0{,}98^{148}=0{,}050$.
\end{solucion}

\begin{solucion}{eje:6-hyperband}
\datos $R=27$, $\eta=3$.
\planteamiento \cref{alg:hyperband}.
\pasos
\begin{enumerate}[label=\arabic*.]
  \item $s_{\max}=\lfloor\log_327\rfloor=3$; $B=4\cdot27=108$.
  \item $s=3$: $n=\lceil4\cdot27/4\rceil=27$, $r=1$: $(27,1)(9,3)(3,9)(1,27)$; recurso $27+27+27+27=108$.
  \item $s=2$: $n=\lceil4\cdot9/3\rceil=12$, $r=3$: $(12,3)(4,9)(1,27)$; $36+36+27=99$.
  \item $s=1$: $n=\lceil4\cdot3/2\rceil=6$, $r=9$: $(6,9)(2,27)$; $54+54=108$.
  \item $s=0$: $n=4$, $r=27$: $(4,27)$; $108$.
  \item Total: $108+99+108+108=423$, equivalente a $423/27=15{,}7$ entrenamientos completos, con $27+12+6+4=49$ configuraciones.
\end{enumerate}
\resultado \textbf{Cuatro brackets; 423 unidades de recurso; 15,7 entrenamientos completos; 49 configuraciones.}
\comprobacion la práctica obtiene $(27,1)(9,3)(3,9)(1,27)$ para el primer bracket, y cada bracket cuesta como mucho $B+(s+1)R\eta^{-s}$ (\cref{prop:bracket}).
\end{solucion}

\begin{solucion}{eje:6-tpe}
\datos $\gamma=0{,}15$; $(\ell,g)$: $(2;0{,}5)$, $(1;1)$, $(0{,}4;0{,}05)$.
\planteamiento \cref{prop:tpe}: $\EI\propto(\gamma+(1-\gamma)g/\ell)^{-1}$.
\pasos
\begin{enumerate}[label=\arabic*.]
  \item $\lambda_1$: $g/\ell=0{,}25$; $1/(0{,}15+0{,}2125)=2{,}759$.
  \item $\lambda_2$: $g/\ell=1$; $1/1=1$.
  \item $\lambda_3$: $g/\ell=0{,}125$; $1/(0{,}15+0{,}1063)=3{,}902$.
\end{enumerate}
\resultado \textbf{Orden $\lambda_3>\lambda_1>\lambda_2$.}
\comprobacion lo que cuenta es el cociente $\ell/g$ ($8$, $4$ y $1$), no $\ell$: $\lambda_3$ es poco frecuente entre las buenas pero casi inexistente entre las malas.
\end{solucion}

\begin{solucion}{eje:6-tpeprueba}
\datos modelo de la \cref{def:tpe}.
\planteamiento repetir la demostración de la \cref{prop:tpe} con detalle.
\pasos
\begin{enumerate}[label=(\alph*)]
  \item $c=\int_{-\infty}^{y^*}(y^*-y)p(y)\,dy$ es la integral de una función no negativa, positiva en un conjunto de probabilidad $\gamma>0$: $c>0$. Que $p(y\mid\lambda)=\ell(\lambda)p(y)/p(\lambda)$ para $y<y^*$ es la regla de Bayes con $p(\lambda\mid y)=\ell(\lambda)$ en esa región.
  \item La demostración solo usa $\Pr[y<y^*]=\gamma$ (para escribir $p(\lambda)=\gamma\ell+(1-\gamma)g$) y que $\ell$ no depende de $y$ por debajo de $y^*$; ambas se cumplen para cualquier distribución continua de $y$ si $y^*$ es su cuantil $\gamma$.
  \item Con $\gamma\to0$, $\EI\propto(g/\ell)^{-1}$ con muy pocos puntos «buenos»: $\ell$ se estima con uno o dos datos y el método explota agresivamente alrededor del mejor; con $\gamma$ grande, $\ell$ se parece a la distribución global y el método explora.
\end{enumerate}
\resultado \textbf{$c>0$; el resultado vale para cualquier distribución continua; $\gamma$ regula exploración y explotación.}
\comprobacion con $g=\ell$ el cociente es 1 en todo $\Lambda$ y la EI es constante: sin información, todas las configuraciones valen lo mismo.
\end{solucion}

\begin{solucion}{eje:6-loo}
\datos $x=(0,1)$, $y=(1,0)$, $\Ky=\bigl(\begin{smallmatrix}1{,}01&0{,}6065\\0{,}6065&1{,}01\end{smallmatrix}\bigr)$.
\planteamiento \cref{teo:loo}.
\pasos
\begin{enumerate}[label=\arabic*.]
  \item $\det\Ky=1{,}0201-0{,}3679=0{,}6522$; $\Ky^{-1}=\bigl(\begin{smallmatrix}1{,}5486&-0{,}9299\\-0{,}9299&1{,}5486\end{smallmatrix}\bigr)$.
  \item $\Ky^{-1}y=(1{,}5486;-0{,}9299)$.
  \item $y_1$: $\mu_{-1}=1-1{,}5486/1{,}5486=0$; $\sigma_{-1}^2=1/1{,}5486=0{,}6458$.
  \item $y_2$: $\mu_{-2}=0+0{,}9299/1{,}5486=0{,}6005$; $\sigma_{-2}^2=0{,}6458$.
  \item Directo para $y_1$ dado $y_2=0$: media $\frac{0{,}6065}{1{,}01}\cdot0=0$; varianza $1{,}01-0{,}6065^2/1{,}01=1{,}01-0{,}3642=0{,}6458$.
\end{enumerate}
\resultado \textbf{$(\mu_{-1},\sigma_{-1}^2)=(0;\ 0{,}6458)$ y $(\mu_{-2},\sigma_{-2}^2)=(0{,}6005;\ 0{,}6458)$.}
\comprobacion $\mu_{-2}=0{,}6005$ coincide con la predicción en $x=1$ con un dato y ruido del \cref{eje:3-gp1} ($0{,}6065/1{,}01$), y la varianza, con aquella más el ruido: $0{,}6358+0{,}01=0{,}6458$.
\end{solucion}

\begin{solucion}{eje:6-parada}
\datos $\max\EI=0{,}004$; mejor $0{,}785$; $\UCB(\hat x)=0{,}795$; $\min\LCB=0{,}770$; $10^6$ piezas; 300~€ por ensayo; ruido $0{,}0126$.
\planteamiento reglas de la §6.6 y \cref{prop:cota}.
\pasos
\begin{enumerate}[label=(\alph*)]
  \item Umbral de Jones: $0{,}01\cdot0{,}785=0{,}00785>0{,}004$: la regla dice parar.
  \item Cota: $0{,}795-0{,}770=0{,}025$~€ por pieza.
  \item Valor esperado de un ensayo más: $0{,}004\cdot10^6=4000$~€ $>300$~€: compensa seguir.
  \item La cota ($0{,}025$) es el doble del ruido de un lote ($0{,}0126$): la incertidumbre del modelo aún domina sobre la de medida, y según el criterio de Makarova \emph{et al.} todavía no se ha llegado al punto de ajustar el ruido.
\end{enumerate}
\resultado \textbf{Jones: parar; cota de regret 0,025~€; criterio económico y cota frente al ruido: seguir.}
\comprobacion el desacuerdo ilustra la observación del texto: la regla del 1\,\% ignora el volumen de producción y el ruido.
\end{solucion}

\begin{solucion}{eje:6-fuga}
\datos 200 configuraciones; se publica el error de validación de la mejor.
\planteamiento \cref{prop:ganador}: el mínimo de estimaciones ruidosas es sesgado hacia abajo.
\pasos
\begin{enumerate}[label=\arabic*.]
  \item Cada error de validación es el verdadero más un ruido (muestra finita de validación, semilla). Al elegir el mínimo de 200, se elige preferentemente una configuración con ruido favorable: $\E[\min]\le\min\E$.
  \item Validación anidada: se separa un conjunto de prueba que no interviene en la búsqueda; dentro del resto se hace la búsqueda con validación (o validación cruzada interna); la configuración elegida se reentrena y se evalúa una sola vez en la prueba. Para estimar la variabilidad, el esquema se repite con varias particiones externas.
\end{enumerate}
\resultado \textbf{El error publicado es optimista; la validación anidada da una estimación insesgada del rendimiento del procedimiento completo.}
\comprobacion en el \cref{ej:hpo}, el error de prueba supera al de validación en 0,05–0,08 en los tres métodos.
\end{solucion}

\begin{solucion}{eje:6-llm}
\datos 5 hiperparámetros; 200~h de GPU; 4~h por entrenamiento completo.
\planteamiento combinar los métodos del tema con el presupuesto: 50 entrenamientos completos equivalentes.
\pasos
\begin{enumerate}[label=\arabic*.]
  \item Espacio: tasa de aprendizaje en $[10^{-5},10^{-3}]$ logarítmica; rango de LoRA categórico ordinal; épocas como \emph{recurso}, no como hiperparámetro; calentamiento en $[0;0{,}1]$; \emph{dropout} en $[0;0{,}2]$.
  \item Método: ASHA (o el \texttt{HyperbandPruner} de Optuna) con muestreo TPE, porque las fidelidades son naturales (fracciones de época), el espacio es mixto y hay posibles trabajadores en paralelo; la BO con PG se reservaría para un ajuste fino final de 2–3 variables.
  \item Recurso: mínimo $1/9$ de época, máximo 5 épocas, $\eta=3$; un ciclo de Hyperband con estos valores cuesta del orden de 15–20 entrenamientos completos, lo que permite 2 ciclos y deja margen.
  \item Semillas: 2 semillas para las 3–5 mejores configuraciones finales (unas 40~h).
  \item Validación: conjunto de validación estratificado por tipo de incidencia para la búsqueda y uno de prueba aparte, que solo se usa al final.
  \item Parada: presupuesto, o cuando la diferencia entre las mejores configuraciones sea menor que la desviación entre semillas.
\end{enumerate}
\resultado \textbf{ASHA con TPE sobre fracciones de época, dos ciclos, reevaluación con semillas y prueba aparte.}
\comprobacion el presupuesto cuadra: unos 35–40 entrenamientos equivalentes para la búsqueda y unos 10 para semillas y prueba, dentro de los 50 disponibles.
\end{solucion}
